\section{Conclusion}
\label{sec:conclusion}

We have shown that commonly used kinase inhibitor selectivity definitions
cluster into two families that measure fundamentally different properties of
binding profiles. We also showed that rank instability within and between families is
mechanistically predictable from binding profile shape, and that three
structurally distinct sources of instability can be identified and
characterized.
These findings motivate four required properties that any well-formed selectivity measure
should satisfy: a reliability threshold
condition, bounded gap sensitivity, distributional consistency, and monotonicity
under weak off-target addition.

No existing definition satisfies all four simultaneously, but we show that a
gated, smooth-hinge effective-target number does. We demonstrate that it recovers its
full-panel ranking from roughly a third of a typical panel, and that it is
part of a whole family of candidate definitions satisfying the four properties.
However, our featured candidate measure is close to identical to the
entropy measure in terms of performance on the given datasets, and
looking for a significantly better measure within the candidate family is future work,
likely requiring additional datasets.
Establishing additional properties that characterize a selectivity measure
uniquely, as Shannon's axioms do for entropy, is a natural next step.

The immediate practical implication is that selectivity assessments should
report the sensitivity of rankings to parameter choice, and should explicitly
flag compounds with no binding above the detection threshold as unreliable.
The longer-term implication is that the field needs a principled basis for
choosing among selectivity definitions.
The need for a principled measure is not specific to selectivity. 
As future and ongoing work, we are developing a general
measurement-theoretic framework for defining and comparing the competing
measures of one concept, reporting where they agree and which pairs they
leave undecided, with kinase selectivity as one instance alongside cases
as varied as biological age and variant effect prediction.

